\documentclass[11pt, fleqn]{amsart}
\usepackage{amsmath, amssymb, amsthm}
\usepackage{hyperref}
\usepackage{mathtools}
\usepackage{listings}
\usepackage{xcolor}
\lstset{
  basicstyle=\ttfamily\small,
  breaklines=true,
  frame=single,
  language=Python,
  keywordstyle=\color{blue},
  commentstyle=\color{gray},
  stringstyle=\color{red}
}
\newtheorem{theorem}{Theorem}[section]
\newtheorem{lemma}[theorem]{Lemma}
\newtheorem{corollary}[theorem]{Corollary}
\newtheorem{proposition}[theorem]{Proposition}
\theoremstyle{remark}
\newtheorem{remark}{Remark}

\title{Spectral graph theory and the matrices $J_n(\overline{j})$, $H_n(\overline{h})$:
remarks on the digraph viewpoint}
\date{12 June 2026}

\begin{document}
\maketitle

\noindent
These remarks concern the two matrices defined after Lemma 2.1 of
\emph{Oscillating spectra associated to certain sequences}, regarded as
weighted adjacency matrices of directed graphs. The short answer is yes:
spectral graph theory says quite a bit, and the directed-graph viewpoint
fits these matrices unusually well. The remarks are organized from
``structural'' to ``potentially useful for the $\mu_n$ project.''

\section{The digraph of $J_n$ has a very rigid cycle structure}

Read $J_n(\overline{j})$ as the weighted adjacency matrix of a digraph on
vertices $1,\dots,n$, with the convention that the entry $A_{uv}$ is the
weight of the edge $u\to v$. The superdiagonal gives a single
``ascending'' path $i \to i+1$ with weight $-i$, and the lower-triangular
Toeplitz part gives ``descending'' edges $i \to k$ ($k \le i$) with weight
$j_{i-k+1}$, including loops of weight $j_1$. Because the \emph{only} way
to ascend is one step at a time, every directed cycle is supported on an
interval $[k,m]$: it climbs $k \to k+1 \to \cdots \to m$ and returns via
the single edge $m \to k$. Its weight is
\begin{equation}\label{eq:cycleweight}
w([k,m]) \;=\; j_{m-k+1}\,(-1)^{m-k}\,\frac{(m-1)!}{(k-1)!}.
\end{equation}

\section{The coefficient theorem for digraphs reproves Theorem 2.3
combinatorially}

The Cvetkovi\'c--Sachs coefficient theorem, which is valid for weighted
digraphs, says that the coefficient of $x^{n-r}$ in $\chi(J_n)$ is
\[
\sum_{L} (-1)^{c(L)}\, w(L),
\]
the sum running over collections $L$ of vertex-disjoint cycles covering
$r$ vertices, where $c(L)$ is the number of cycles in $L$ and $w(L)$ is
the product of the edge weights occurring in $L$. By the previous
section, those collections are exactly sets of pairwise disjoint
intervals in $\{1,\dots,n\}$ --- that is, the coefficients are signed
sums over partial composition-like structures, which is morally why
partition and composition identities (MacDonald, Vein--Dale) appear at
all in this setting.

This was verified numerically (see the script in
Section~\ref{sec:script}): for a random integer sequence $\overline{j}$
with $j_1=1$ and $n=7$, the cycle-cover sums match both \texttt{numpy}'s
characteristic polynomial and the formula of Theorem 2.3,
\[
\binom{n}{r}\, r!\, h_r\, (-1)^r,
\]
exactly.

A corollary worth noting: for digraphs the spectrum depends \emph{only}
on the multiset of cycle products. (Diagonal similarity is a ``gauge
transformation'' that preserves cycle weights, and the coefficient
theorem shows the characteristic polynomial is determined by them.)
Hence the interval weights \eqref{eq:cycleweight} are the complete
intrinsic data behind the spectra under study. Any answer to the
question ``what is it about $\chi(J_n)$ that induces the observed
oscillations?'' must be expressible in terms of them.

\section{Strong connectivity and its standard dividends}

The digraph is strongly connected (one can ascend stepwise and descend
in one hop), so $J_n$ is irreducible; and since the superdiagonal is
nowhere zero, $J_n$ is an unreduced Hessenberg matrix, hence
nonderogatory --- every eigenvalue has geometric multiplicity $1$, and
$\chi(J_n)$ coincides with the minimal polynomial.

Irreducibility upgrades the Ger\v{s}gorin disc theorem via Taussky's
theorem: an eigenvalue lying on the boundary of the union of the discs
must lie on the boundary of \emph{every} disc. More interestingly,
Brualdi's cycle-based refinement of Ger\v{s}gorin is spectral graph
theory tailor-made for the present situation: it localizes eigenvalues
using products of deleted row sums taken \emph{around directed cycles}
of the digraph, and here the cycles are exactly the intervals described
above.

\section{The Ger\v{s}gorin picture explains, geometrically, what
deformation does}

Every diagonal entry of $J_n$ equals $j_1$, so all Ger\v{s}gorin discs
share one center. Deformation by $c$ replaces every loop weight by $c$
(and shifts the descending-edge weights), so it translates the common
disc center to $c$ while reorganizing the radii. In graph language,
deformation is a uniform reweighting of the loops --- the spectrum is
pinned to circulate around $c$. This gives at least a heuristic for why
the deformed spectra look more regular: the eigenvalues live in an
annulus-like region about a fixed center, and $\mu_n$ measures the
distance from $0$ to a cloud orbiting $c$.

\section{A concrete lever for bounding $\mu_n$ away from zero}

Since
\[
\prod_i |\lambda_i| \;=\; |\det J_n| \;=\; n!\,|h_n|,
\]
any upper bound $\rho$ on the spectral radius coming from the graph
(Ger\v{s}gorin row sums in terms of the $j$'s, or Brualdi's cycle
bounds) yields
\[
\mu_n \;\ge\; \frac{n!\,|h_n|}{\rho^{\,n-1}}.
\]
This is exactly the shape of statement the project wants --- it reduces
``the spectrum is zero-free'' to controlling the top of the spectrum
graph-theoretically, plus knowing $h_n \neq 0$. For the Lehmer
application the logic unfortunately runs the wrong way ($h_n \neq 0$ is
the desired conclusion, not a hypothesis). Still, for the deformed
matrices, where only zero-freeness is needed, it is potentially usable.

\section{Caveats, and the matrix $H_n$}

One important caveat: the classical heavy machinery of spectral graph
theory (Perron--Frobenius, eigenvalue interlacing, expander-type
bounds) assumes nonnegative or symmetric matrices. The weights here are
signed and the matrices are far from normal, so the spectra are
genuinely complex point sets and those tools do not apply. What
survives is exactly the digraph toolkit above: the coefficient theorem,
irreducibility, and cycle-based inclusion regions.

Everything said for $J_n$ transfers to $H_n(\overline{h})$ with the
obvious changes: the superdiagonal weights are all $+1$, and the
first-column edges $i \to 1$ carry weight $i\,h_i$. Consequently the
cycles of the $H_n$ digraph are the loops, the intervals $[k,m]$ with
$k \ge 2$, and the intervals $[1,m]$, the last carrying the extra
factor $m$.

\section{The Python verification script}\label{sec:script}

The following script (i) builds $J_n(\overline{j})$ for a random
integer sequence with $j_1 = 1$; (ii) computes its characteristic
polynomial with \texttt{numpy}; (iii) recomputes the coefficients from
the digraph coefficient theorem, enumerating all sets of pairwise
disjoint intervals and using the cycle weights
\eqref{eq:cycleweight}; and (iv) compares both against the formula of
Theorem 2.3, with $h$ recovered from $j$ via equation (D1). All three
computations agree.

\begin{lstlisting}
import numpy as np
import random
from math import comb, factorial

random.seed(1)
n = 7
# random integer j-sequence with j_1 = 1 (as in Lemma 2.2)
j = [None, 1] + [random.randint(-4, 4) for _ in range(2, n + 1)]  # j[1..n]

# Build J_n(jbar): lower part Toeplitz j_{i-k+1}, superdiagonal -i
J = np.zeros((n, n))
for i in range(1, n + 1):
    for k in range(1, n + 1):
        if k <= i:
            J[i - 1, k - 1] = j[i - k + 1]
        elif k == i + 1:
            J[i - 1, k - 1] = -i

# Characteristic polynomial via numpy:
# coefficients of x^n + a_1 x^{n-1} + ...
cp_numeric = np.poly(J)

# --- digraph coefficient theorem ---
# Cycles of the digraph = intervals [k,m]:
#   weight = j_{m-k+1} * (-1)^(m-k) * (m-1)!/(k-1)!
def cycle_weight(k, m):
    return j[m - k + 1] * ((-1) ** (m - k)) * factorial(m - 1) // factorial(k - 1)

# Coefficient of x^{n-r}: sum over sets of disjoint intervals
# covering r vertices, with sign (-1)^{#cycles}  (det(xI - A) convention)
intervals = [(k, m) for k in range(1, n + 1) for m in range(k, n + 1)]

def disjoint_sets(ivs):
    """Enumerate all sets of pairwise disjoint intervals."""
    out = [[]]
    def rec(start, chosen):
        for idx in range(start, len(ivs)):
            iv = ivs[idx]
            if all(iv[0] > c[1] or iv[1] < c[0] for c in chosen):
                chosen.append(iv)
                out.append(list(chosen))
                rec(idx + 1, chosen)
                chosen.pop()
    rec(0, [])
    return out

coeffs = [0] * (n + 1)  # coeffs[r] multiplies x^{n-r}
coeffs[0] = 1
for S in disjoint_sets(intervals):
    if not S:
        continue
    r = sum(m - k + 1 for k, m in S)
    w = 1
    for k, m in S:
        w *= cycle_weight(k, m)
    coeffs[r] += ((-1) ** len(S)) * w

print("numpy char poly coeffs:    ", np.round(cp_numeric, 6))
print("digraph cycle-cover coeffs:", coeffs)
print("match:", np.allclose(cp_numeric, coeffs))

# --- confirm these equal Theorem 2.3's binom(n,r) r! h_r (-1)^r,
#     with h recovered from j via equation (D1) ---
h = [1]
for m in range(1, n + 1):
    s = sum(j[r] * h[m - r] for r in range(1, m + 1))
    h.append(s / m)
thm = [comb(n, r) * factorial(r) * h[r] * (-1) ** r for r in range(n + 1)]
print("Theorem 2.3 coeffs:        ", thm)
\end{lstlisting}

\noindent
Output for the seeded run ($n = 7$):
\begin{lstlisting}[language={}]
numpy char poly coeffs:     [ 1. -7. -21. 385. -805. -609. 6433. -8041.]
digraph cycle-cover coeffs: [1, -7, -21, 385, -805, -609, 6433, -8041]
match: True
Theorem 2.3 coeffs:         [1, -7.0, -21.0, 385.0, -805.0, -609.0,
                             6433.0, -8041.0]
\end{lstlisting}

\end{document}
